% 苏布拉马尼扬·钱德拉塞卡（综述）
% license CCBYSA3
% type Wiki

本文根据 CC-BY-SA 协议转载翻译自维基百科\href{https://en.wikipedia.org/wiki/Subrahmanyan_Chandrasekhar}{相关文章}

苏布拉马尼扬·钱德拉塞卡（Subrahmanyan Chandrasekhar，/ˌtʃəndrəˈʃeɪkər/，音译：昌德拉谢卡；[4] 泰米尔语：சுப்பிரமணியன் சந்திரசேகர்，罗马化：Cuppiramaṇiyaṉ Cantiracēkar；1910年10月19日－1995年8月21日）[5] 是一位印裔美国理论物理学家，在**恒星结构、恒星演化以及黑洞**等领域作出了卓越贡献。他在职业生涯中也曾将相当一部分时间投入于**流体力学研究**，特别是**稳定性与湍流理论**，并在这些领域取得了重要成果。钱德拉塞卡与**威廉·A·福勒（William A. Fowler）**因在“**与恒星结构与演化有关的重要物理过程的理论研究**”方面的贡献，共同获得**1983年诺贝尔物理学奖**。他对**恒星演化的数学处理**为现代天体物理学提供了众多关于**大质量恒星及黑洞晚期演化阶段**的理论模型。[6][7]许多重要的科学概念、机构与发明均以他命名，包括著名的**钱德拉塞卡极限与**钱德拉X射线天文台。[8]
